\chapter{Analyse du besoin et objectifs}
\label{ch:besoin}

\section{Problématique métier}

Un gestionnaire de stock doit pouvoir répondre rapidement aux questions suivantes :
\begin{enumerate}
  \item Quels produits sont en rupture ou proches du seuil minimum ?
  \item Quelles entrées et sorties ont eu lieu, par qui, et pour quel motif ?
  \item Comment organiser le catalogue par familles (catégories) ?
\end{enumerate}

Ces besoins imposent une base relationnelle cohérente, des règles métier transactionnelles (un mouvement de sortie ne doit pas rendre le stock négatif) et une interface qui reflète l'état en temps quasi réel après chaque opération.

\section{Acteurs et rôles}

\begin{table}[htbp]
  \centering
  \caption{Acteurs et permissions}
  \label{tab:acteurs}
  \begin{tabular}{@{}llp{7cm}@{}}
    \toprule
    \textbf{Rôle} & \textbf{Compte démo} & \textbf{Capacités} \\
    \midrule
    Administrateur & \texttt{admin@stock.tn} & CRUD complet, suppression catégories/produits \\
    Opérateur & \texttt{operateur@stock.tn} & Consultation, création mouvements et entités (sauf suppressions réservées admin) \\
    \bottomrule
  \end{tabular}
\end{table}

L'authentification repose sur un jeton JWT transmis dans l'en-tête \texttt{Authorization: Bearer}. Le backend résout l'utilisateur courant à chaque requête protégée via la dépendance \texttt{obtenir\_utilisateur\_courant}.

\section{Exigences fonctionnelles}

\begin{table}[htbp]
  \centering
  \caption{Exigences fonctionnelles principales}
  \label{tab:exigences-fonctionnelles}
  \small
  \begin{tabular}{@{}p{2.2cm}p{11cm}@{}}
    \toprule
    \textbf{ID} & \textbf{Description} \\
    \midrule
    F1 & Connexion par email et mot de passe, émission d'un jeton d'accès. \\
    F2 & CRUD catégories (nom, description optionnelle). \\
    F3 & CRUD produits (nom, référence SKU unique, quantité, seuil d'alerte, catégorie). \\
    F4 & Enregistrement de mouvements \texttt{entree} ou \texttt{sortie} avec mise à jour atomique du stock. \\
    F5 & Rejet d'une sortie si le stock deviendrait négatif (\texttt{stock\_insuffisant}). \\
    F6 & Tableau de bord : compteurs produits, catégories, mouvements, liste des alertes. \\
    F7 & Page alertes : produits dont \texttt{quantite <= seuil\_alerte}. \\
    \bottomrule
  \end{tabular}
\end{table}

\section{Exigences non fonctionnelles}

\begin{table}[htbp]
  \centering
  \caption{Exigences non fonctionnelles}
  \label{tab:exigences-nf}
  \small
  \begin{tabular}{@{}p{2.2cm}p{11cm}@{}}
    \toprule
    \textbf{ID} & \textbf{Description} \\
    \midrule
    NF1 & \textbf{Conteneurisation} : exécution reproductible via Docker. \\
    NF2 & \textbf{Orchestration} : déploiement sur Kubernetes avec Ingress, sondes de santé, ServiceMonitor. \\
    NF3 & \textbf{IaC} : infrastructure déclarative (Terraform + Helm). \\
    NF4 & \textbf{Observabilité} : métriques Prometheus sur \texttt{/metrics}, dashboard Grafana. \\
    NF5 & \textbf{CI/CD} : validation automatique et déploiement production (GitHub Actions). \\
    NF6 & \textbf{Sécurité} : mots de passe hashés (bcrypt), secrets Kubernetes pour données sensibles. \\
    NF7 & \textbf{Documentation opérationnelle} : README, Makefile, scripts \texttt{make demo}. \\
    \bottomrule
  \end{tabular}
\end{table}

\section{Objectifs techniques}

Les objectifs techniques déclinés dans l'implémentation sont :

\begin{enumerate}
  \item \textbf{Séparation des couches} : frontend statique, API stateless, base PostgreSQL.
  \item \textbf{Deux chemins de démonstration} : Compose local (\texttt{make dev}) et minikube (\texttt{make cluster}) avec le même chart Helm qu'en CI.
  \item \textbf{Automatisation} : \texttt{make demo} en local ; CI sur chaque push/PR.
\end{enumerate}

\section{Hors périmètre}

Les éléments suivants sont explicitement hors périmètre de la version documentée :
\begin{itemize}
  \item déploiement production VPS automatique (workflow \texttt{deploy.yml} manuel uniquement) ;
  \item haute disponibilité multi-nœuds, autoscaling horizontal, TLS terminé par certificats réels ;
  \item gestion fine des permissions au-delà de deux rôles ;
  \item application mobile native ou mode hors-ligne.
\end{itemize}
